% ---------------------------------------------------------------
% Code-Blöcke für YAML und Bash (in die Präambel einbinden)
%   \input{code_listings}
% ---------------------------------------------------------------
\usepackage{xcolor}
\usepackage{listings}
\usepackage[most]{tcolorbox}

% Farben
\definecolor{codebg}{HTML}{F6F8FA}
\definecolor{codeframe}{HTML}{D0D7DE}
\definecolor{codetitlebg}{HTML}{EAEEF2}
\definecolor{codetitle}{HTML}{24292F}
\definecolor{codecomment}{HTML}{6E7781}
\definecolor{codekey}{HTML}{0550AE}
\definecolor{codestring}{HTML}{0A3069}
\definecolor{codevalue}{HTML}{953800}
\definecolor{codekeyword}{HTML}{CF222E}

% Gemeinsame Basis (inkl. Umlaute in Listings)
\lstdefinestyle{codebase}{
basicstyle=\ttfamily\small,
columns=flexible,
keepspaces=true,
showstringspaces=false,
breaklines=true,
breakatwhitespace=true,
postbreak=\mbox{\textcolor{codecomment}{$\hookrightarrow$}\space},
tabsize=2,
extendedchars=true,
literate={ä}{{\"a}}1 {ö}{{\"o}}1 {ü}{{\"u}}1
{Ä}{{\"A}}1 {Ö}{{\"O}}1 {Ü}{{\"U}}1 {ß}{{\ss}}1,
}

% YAML: Schlüssel blau, Werte orange, Strings dunkelblau, Kommentare grau
\newcommand\YAMLcolonstyle{\color{codecomment}}
\newcommand\YAMLvaluestyle{\color{codevalue}}
\lstdefinelanguage{yaml}{
    comment=[l]{\#},
    morestring=[b]",
    morestring=[b]',
    moredelim=**[il][\YAMLcolonstyle{:}\YAMLvaluestyle]{:},
    keywords={true,false,null},
    sensitive=false,
}
\lstdefinestyle{yaml}{
    style=codebase,
    language=yaml,
    basicstyle=\ttfamily\small,
    commentstyle=\color{codecomment},
    stringstyle=\color{codestring},
    keywordstyle=\color{codekeyword},
}

% Bash
\lstdefinestyle{bash}{
    style=codebase,
    language=bash,
    basicstyle=\ttfamily\small,
    keywordstyle=\color{codekeyword},
    morekeywords={docker,pull,run,network,create},
    commentstyle=\color{codecomment},
    stringstyle=\color{codestring},
}

% Rahmen für beide Blocktypen; optionales Argument z.B. [title={config.yaml}]
\tcbset{
    codebox/.style={
            enhanced, breakable, listing only,
            colback=codebg, colframe=codeframe,
            colbacktitle=codetitlebg, coltitle=codetitle,
            fonttitle=\small\ttfamily,
            boxrule=0.5pt, arc=3pt,
            left=6pt, right=6pt, top=3pt, bottom=3pt,
            toptitle=2pt, bottomtitle=2pt,
        }
}
% colupper setzt die Grundfarbe (YAML-Schlüssel), auch über Seitenumbrüche hinweg
\newtcblisting{yamlcode}[1][]{codebox, colupper=codekey, listing options={style=yaml}, #1}
\newtcblisting{bashcode}[1][]{codebox, colupper=codetitle, listing options={style=bash}, #1}